\medskip

On se propose d'étudier la concentration dans le sang d'un médicament ingéré par une personne pour la première fois. Soit $t$ le temps (en heures) écoulé depuis l'ingestion de ce médicament.
On admet que la concentration de ce médicament dans le sang, en gramme par litre de sang, est modélisée par une fonction $f$ de la variable $t$ définie sur l'intervalle $[0~;~+ \infty[$.


\textbf{Partie A : lectures graphiques}

\begin{center}
\psset{unit=1.6cm,arrowsize=2pt 3}
\begin{pspicture*}(-0.4,-0.6)(7.6,2.4)
\psgrid[unit=0.8cm,gridlabels=0pt,gridcolor=gray!50,subgridcolor=lightgray!50](-1,-2)(16,5)
\psaxes[linewidth=1.25pt,labelFontSize=\scriptstyle]{->}(0,0)(-0.4,-0.6)(7.6,2.4)
\psplot[plotpoints=2000,linewidth=1.25pt,linecolor=blue]{0}{7.5}{5 x mul 2.71828 x exp div}
\uput[d](6.8,-.2){\footnotesize temps en heures}
\uput[r](0,2.2){\footnotesize concentration en g/L}

\psset{linecolor=red,linestyle=dashed,linewidth=1.2pt}
\psline(0,1)(7.5,1)
\psline(0.25,0)(0.25,1) \psline(2.5,0)(2.5,1)
\psline[linestyle=solid,linewidth=2pt]{[-]}(0.25,0)(2.5,0)
\psline[linecolor=blue](1,0)(1,1.84)
\end{pspicture*}
\end{center}

On a représenté ci-dessus la courbe représentative de la fonction $f$. 

Avec la précision permise par le graphique, on donne sans justification :



\begin{enumerate}
\item Le temps écoulé depuis l'instant de l'ingestion de ce médicament et l'instant où la concentration de médicament dans le sang est maximale selon ce modèle: 1 heure.

\item L'ensemble des solutions de l'inéquation $f(t) \geqslant 1$: l'intervalle $\left [0,25\;;\; 2,5\strut\right ]$.

\item La convexité de la fonction $f$ sur l' intervalle [0~;~8]: la fonction semble concave entre 0 et 2, puis convexe.
\end{enumerate}

\medskip

\textbf{Partie B : détermination de la fonction} \boldmath $f$ \unboldmath

\medskip

On considère l'équation différentielle  $(E)$: $y' +y = 5\e^{-t}$,

d'inconnue $y$, où $y$ est une fonction définie et dérivable sur l'intervalle $[0~;~+ \infty[$.

On admet que la fonction $f$ est une solution de l'équation différentielle $(E)$.



\begin{enumerate}
\item On résout l'équation différentielle $(E')$: $y' + y = 0$.

D'après le cours, on sait que les équations différentielles de la forme $ay'+by=0$ ont des solutions $y$ s'écrivant $y(t) = k\e^{-\frac{b}{a}t}$ où $k$ est un réel quelconque.

Donc l'équation $(E')$ a pour solutions les fonctions $y$ s'écrivant $y(t)=k\e^{-t}$ où $k\in\R$.

\item Soit $u$ la fonction définie sur l'intervalle $[0~;~ +\infty[$ par $u(t) = at\e^{-t}$ avec $a \in \R$.



La fonction $u$ est solution de l'équation différentielle $(E)$ si et seulement si $u'(t)+u(t)=5\e^{-t}$.

$u(t)=at\e^{-t}$ donc $u'(t)= a\times \e^{-t} + at\times (-1)\e^{-t}= a\e^{-t}-at\e^{-t}$

$u'(t)+u(t)=5\e^{-t}
\iff a\e^{-t}-at\e^{-t} + at\e^{-t}=5\e^{-t}
\iff a\e^{-t}=5\e^{-t}
\iff a=5$ car $\e^{-t}\neq 0$ pour tout $t$.

Donc la fonction $u$ définie par $u(t)=5t\e^{-t}$ est une solution de $(E)$.

\item
La solution générale de l'équation $(E)$ est la somme d'une solution particulière de $(E)$ et de la solution générale de l'équation sans second membre $(E')$: ce sont donc les fonctions $f$ définies par $f(t)=k\e^{-t}+5t\e^{-t}$ où $k\in\R$.

\item La personne n'ayant pas pris ce médicament auparavant, on admet que $f(0) = 0$.



$f(0)=0 \iff k\e^{0}+5\times 0 \times \e^{0}=0 \iff = 0$

L'expression de la fonction $f$ est $f(t)=5t\e^{-t}$.

\end{enumerate}

\medskip

\textbf{Partie C : étude de la fonction} \boldmath $f$ \unboldmath

\medskip

Dans cette partie, on admet que $f$ est définie sur l'intervalle $[0~;~+\infty[$ par $f(t) = 5t \e^{-t}$.



\begin{enumerate}
\item 
On sait que $\ds\lim_{t\to +\infty} \dfrac{\e^{t}}{t}=+\infty$ donc $\ds\lim_{t\to +\infty} \dfrac{t}{\e^{t}}=0$, donc $\ds\lim_{t\to +\infty} 5t\e^{-t}=0$, et donc $\ds\lim_{t\to +\infty} f(t)=0$

Cela signifie que la concentration en médicament devient nulle quand le temps augmente indéfiniment.



\item 
$f$ est le produit de fonctions dérivables sur $\R$ donc $f$ est dérivable sur $[0~;~+\infty[$.

$f(t)=5t\e^{-t}$ donc $f'(t) = 5\times \e^{-t} + 5t\times (-1)\e^{-t} = \left (5-5t\right )\e^{-t}$

On étudie le signe de $f'(t)$ sur $[0~;~+\infty[$.

\[ \begin{tablvar}[intervalwidth=6em,stretch=1.2]{2}
 \hline
t & 0 & & 1 & & +\infty\\
 \hline
5-5t & & + & \barre[0] & - & \\
 \hline
\e^{-t} & & + & \barre & + & \\
 \hline
f'(t) & & + & \barre[0] & - & \\
 \hline
 \end{tablvar}\]

$f(0)=0$ et $f(1)= 5\times 1 \times \e^{-1} = 5\e^{-1}$

On établit le tableau de variation complet de $f$ sur  $[0~;~+\infty[$.

\[\begin{tablvar}[6em]{2}
\hline
t & 0 && 1 && +\infty \\
\hline
f'(t) & & + & \barre[0] & - & \\
\hline
\variations{\mil{f} & \bas{0} && \haut{5\e^{-1}} &&  \bas{0}} 
\hline
\end{tablvar}\]

\item On veut montrer qu'il existe deux réels $t_1$ et $t_2$ tels que $f\left(t_1\right) = f\left(t_2\right) = 1$.

Le maximum de la fonction $f$ est $f(1)=5\e^{-1}\approx 1,84>1$.

On complète le tableau de variation de $f$.

\[\begin{tablvar}[6em]{2}
\hline
t & 0 & \vr{t_1} & 1 &  \vr{t_2} & +\infty\\
\hline
\variations{\mil{f} & \bas{0} & \vr{1} & \haut{5\e^{-1}} & \vr{1} & \bas{0}} 
\hline
\end{tablvar}\]

\begin{list}{\textbullet}{}
\item Sur l'intervalle $]0~;~1[$, la fonction $f$ est continue et strictement croissante. De plus $1 \in\, \left ]f(0)~;~f(1)\right [$. Donc, d'après le corollaire du théorème des valeurs intermédiaires, l'équation $f(t)=1$ admet une solution unique dans $]0~;~1[$. On l'appelle $t_1$.

D'après la calculatrice: $f(0,25) \approx 0,97  < 1$ et $f(0,26) \approx 1,002 >1$ donc $0,25$ est une valeur approchée de $t_1$ à $10^{-2}$.

\item Sur l'intervalle $]1~;~+\infty[$, la fonction $f$ est continue et strictement décroissante. De plus $1 \in\, \left ]\ds\lim_{t\to +\infty} f(t)~;~f(1)\right [$. Donc, d'après le corollaire du théorème des valeurs intermédiaires, l'équation $f(t)=1$ admet une solution unique dans $]1~;~+\infty[$. On l'appelle $t_2$.

D'après la calculatrice: $f(2,54) \approx  1,002 > 1$ et $f(2,55) \approx 0,996  <1$ donc $2,54$ est une valeur approchée de $t_2$ à $10^{-2}$.
\end{list}



\item Pour une concentration du médicament supérieure ou égale à 1 gramme par litre de sang, il y a un risque de somnolence. C'est donc quand $f(t) \geqslant 1$.

$f(t)\geqslant 1$ pour $t\in  \left [ t_1~;~t_2\strut \right ]$, soit pour une durée égale à $t_2-t_1$ qui vaut environ $2,54 - 0,25$ soit $2,29$ heures; cela fait environ 2 heures et 17 minutes.


\end{enumerate}

\medskip

\textbf{Partie D : concentration moyenne}

\medskip

La concentration moyenne du médicament (en gramme par litre de sang) durant la première heure est donnée par :
$T_m = \displaystyle\int_0^1 f(t) \d t$
où $f$ est définie sur $[0~;~+\infty[$ par $f(t) = 5t\e^{-t}$.



On utilise une intégration par parties:
$\ds\int_{a}^{b} u(t)v'(t)\d t =\left [ u(t)v(t)\strut \right ]_{a}^{b} -\int_{a}^{b} u'(t) v(t) \d t$.

On pose
$\left \{
\begin{array}{l}
u(t)=5t\\
v'(t)=\e^{-t}
\end{array}
\right .$
donc
$\left \{
\begin{array}{l}
u'(t)=5\\
v(t)= - \e^{-t}
\end{array}
\right .$

$\begin{aligned}[t]
T_m
& = \displaystyle\int_0^1 5t\e^{-t} \d t
 = \left [5t\times \left (-\e^{-t} \right )\right ]_{0}^{1}- \ds\int_{0}^{1} 5\times \left ( -\e^{-t} \right ) \d t
  = \left [-5 t\e^{-t} \right ]_{0}^{1}- 5\left [ \e^{-t} \right ]_{0}^{1}\\
& = \left (-5\e^{-1}-0\right )-\left (-5\e^{-1} + 5\e^{0} \right )
= 5-10\e^{-1} \approx 1,32
\end{aligned}$



\bigskip
